\documentclass[12 pt, letterpaper]{hmcpset}
\usepackage{hw-preamble}

\name{Jayden Thadani}
\class{MATH 171 --- Abstract Algebra I}
\assignment{Homework 14}
\duedate{12th April 2024}

\begin{document}

\begin{problem}[Dummit \& Foote \S8.1 \#3]
	Let $R$ be a Euclidean domain. Let $m$ be the minimum integer in the set of norms of nonzero elements of $R$. Prove that every nonzero element of $R$ of norm $m$ is a unit. Deduce that a nonzero element of norm zero (if such an element exists) is a unit.
\end{problem}

\begin{solution}
	Since $R$ is a Euclidean domain, for any $x \in R$ with $N(x) = m$, we have $1 = qx + r$ where $N(r) < N(x) = m$ or $r = 0$. Since $N(r) < m$ would be a contradiction, we must have $r = 0$, and thus, $qx = 1$, giving us that $x$ is a unit.
\end{solution}

\pagebreak

\begin{problem}[Dummit \& Foote \S8.1 \#4(a)]
	Let $R$ be a Euclidean domain. Prove that if $(a, b) = 1$ and $a$ divides $bc$, then $a$ divides $c$. More generally, show that if $a$ divides $bc$ with nonzero $a, b$, then $\frac{a}{(a, b)}$ divides $c$.
\end{problem}

\begin{solution}
	If $a$ divides $bc$, then we have $bc = a q$ for some $q \in R$. However, since $d = \gcd (a, b) \in (a, b)$ we must have $d = ax + by$ for some $x, y \in R$.
	
	Since we have that $a/d$ and $b/d$ are both in $R$, we can write $1 = (a/d)x + (b/d)y$. Then by multiplying by $c$ on both sides and using the fact that $bc = aq$, we get
	\[
		\begin{split}
			c & = \frac{a}{d} cx + \frac{b}{d} cy \\
			  & = \frac{a}{d} cx + \frac{b}{d} qy \\
			  & = \frac{a}{d} (cx + qy).
		\end{split}
	\]
	and thus, $(a/d) \mid c$ as desired.
\end{solution}

\pagebreak

\begin{problem}[Dummit \& Foote \S8.2 \#2]
	Prove that any two nonzero elements of a PID have a least common multiple.
\end{problem}

\begin{solution}
	Let $R$ be a PID. Then, for any nonzero $a, b \in R$, consider $c \in R$ such that $(c) = (a) \cap (b)$. 

	First we will show that $c$ is a multiple of $a$ and $b$. We know that $(c) = (a) \cap (b) \subset (a)$ and $(c) = (a) \cap (b) \subset (b)$ and thus, we know that $c$ is a multiple of both $a$ and $b$.

	Suppose that $d$ is some multiple of both $a$ and $b$. Then we have $(d) \subset (a)$ and $(d) \subset (b)$. Thus, $(d) \subset (a) \cap (b)$, giving us $(d) \subset (c)$ and thus, $c \mid d$. That is, $c$ is the least common multiple of $a$ and $b$.
\end{solution}

\pagebreak

\begin{problem}[Dummit \& Foote \S8.2 \#3]
	Prove that the quotient of a PID by a prime ideal is again a PID.
\end{problem}

\begin{solution}
	Let $R$ be a PID and let $P$ be a prime ideal in $R$. Since $P$ is prime, we know that $R/P$ is an integral domain. We just need to show that every ideal is principal. 

	By the fourth isomorphism theorem, we know that any ideal $J \subset R/P$ must be of the form $J = I/P$ where $I \subset R$ is an ideal such that $P \subset I$. Since $R$ is a PID, $I = (i)$ for some $i \in R$. That is, every $x \in I$ is $x = ir$ for some $r \in R$. Then, every $x + P \in I/P$, is $(ir) + P$ for some $r \in R$ and thus, $x + P = (i + P)(r + P)$ for some $r + P \in I/P$. Thus, any ideal $J \subset R/P$ is just generated by $i + R$. That is, $J = (i + R)$ where $J = (i)/P$.
\end{solution}

\end{document}
